\newpage
\section{Analiza problemu}		%2
%Napisać gdzie używa się tego algorytmu
%Opisać sposób działania programu/algorytmu
%Napisać sposób wykorzystania algorytmu po przez wykonanie przykładu (np. mnożenie macierzy - wykonać ręcznie przykład z mnożeniem macierzy pokazujący jak mnoży się macierz ręcznie)
%Jeśli zadanie zakłada przedstawienie jakiegoś narzędzia (np. git, AI) należy opisać narzędzie

\hspace{0.60cm}
W niniejszym rozdziale dokonano szczegółowej analizy teoretycznej oraz praktycznej struktury danych w postaci drzewa poszukiwań binarnych (ang. \textit{Binary Search Tree} -- BST). Omówiono praktyczne obszary zastosowań tego algorytmu w nowoczesnej informatyce, formalny sposób działania poszczególnych operacji, jak również przedstawiono pełny przykład wykonania operacji na drzewie krok po kroku. Ponadto opisano narzędzia inżynierii oprogramowania wykorzystane w trakcie realizacji projektu, w tym system kontroli wersji Git, platformę GitHub oraz generator dokumentacji Doxygen.

\subsection{Zastosowanie drzew poszukiwań binarnych w informatyce}

\hspace{0.60cm}
Drzewo poszukiwań binarnych jest jedną z najbardziej fundamentalnych dynamicznych struktur danych. Ze względu na swoje właściwości wyszukiwania oraz możliwości szybkiej modyfikacji, algorytmy oparte na drzewach BST znajdują szerokie zastosowanie w wielu gałęziach informatyki:

\begin{enumerate}
    \item \textbf{Systemy zarządzania bazami danych (DBMS) i indeksowanie:}
    Struktury wariantowe drzew BST (takie jak drzewa AVL, drzewa czerwono-czarne czy B-drzewa) stanowią podstawę indeksowania relacyjnych baz danych. Pozwalają na natychmiastowe zlokalizowanie rekordu o określonym kluczu w czasie logarytmicznym $O(\log n)$ bez konieczności liniowego przeglądania całej bazy.
    
    \item \textbf{Implementacja słowników i zbiorów w bibliotekach standardowych:}
    Wiele języków programowania wykorzystuje zrównoważone drzewa BST do realizacji kontenerów asocjacyjnych. Przykładem jest biblioteka STL w języku C++, gdzie kontenery takie jak \texttt{std::map} oraz \texttt{std::set} są domyślnie zaimplementowane w oparciu o drzewa czerwono-czarne.
    
    \item \textbf{Kompresja danych (Drzewa Huffmana):}
    W algorytmach bezstratnej kompresji danych, takich jak kodowanie Huffmana, dynamiczne struktury drzewiaste są wykorzystywane do budowania prefiksowych kodów znaków o zmiennej długości na podstawie częstotliwości ich występowania.
    
    \item \textbf{Drzewa wyrażeń i parsowanie składniowe (AST):}
    W kompilatorach i interpreterach języków programowania drzewa są wykorzystywane jako abstrakcyjne drzewa składniowe (ang. \textit{Abstract Syntax Tree} -- AST) do reprezentowania struktury gramatycznej kodu źródłowego oraz ewaluacji wyrażeń arytmetycznych i logicznych.
    
    \item \textbf{Grafika komputerowa i podział przestrzeni (BSP):}
    Drzewa binarnego podziału przestrzeni (ang. \textit{Binary Space Partitioning} -- BSP) pomagają w sortowaniu i wyświetlaniu obiektów trójwymiarowych w kolejności od przodu do tyłu lub odwrotnie, eliminując potrzebę wykonywania kosztownych obliczeń widoczności powierzchni.
\end{enumerate}

\subsection{Sposób działania algorytmu i podstawowe operacje}

\hspace{0.60cm}
Struktura BST składa się z węzłów połączonych wskaźnikami. Każdy węzeł zawiera przechowywaną wartość (\texttt{data}) oraz dwa wskaźniki poddrzew: lewego (\texttt{left}) i prawego (\texttt{right}). Nadrzędną zasadą porządkującą strukturę jest relacja:
\begin{equation}
    \forall_{v \in \text{Left}(N)} \text{val}(v) < \text{val}(N) \quad \text{oraz} \quad \forall_{u \in \text{Right}(N)} \text{val}(u) \ge \text{val}(N)
\end{equation}

\hspace{0.60cm}
Poniżej opisano mechanizm działania kluczowych algorytmów zaimplementowanych w klasie \texttt{BSTTree}:

\begin{itemize}
    \item \textbf{Dodawanie elementu (Insertion):}
    Rozpoczyna się od korzenia. Jeśli drzewo jest puste, nowy węzeł staje się korzeniem. W przeciwnym razie porównuje się wstawianą wartość z wartością w aktualnym węźle. Jeśli jest mniejsza, przechodzi się rekurencyjnie do lewego poddrzewa; jeśli jest większa lub równa -- do prawego. Proces powtarza się do momentu natrafienia na wskaźnik \texttt{nullptr}, w miejscu którego tworzony jest nowy węzeł.
    
    \item \textbf{Wyszukiwanie drogi do elementu (Search / Path finding):}
    Algorytm rozpoczyna przeszukiwanie od korzenia, dopisując każdy odwiedzany węzeł do wektora ścieżki. Przechodzenie w dół drzewa odbywa się zgodnie z relacją porządku. Wyszukiwanie kończy się powodzeniem po odnalezieniu szukanego klucza lub niepowodzeniem po osiągnięciu końca gałęzi (\texttt{nullptr}).
    
    \item \textbf{Algorytmy przechodzenia drzewa (Traversals):}
    \begin{itemize}
        \item \textit{Preorder (korzeń, lewe, prawe):} Służy do kopiowania lub zapisu struktury drzewa.
        \item \textit{Inorder (lewe, korzeń, prawe):} Zwraca uporządkowaną (rosnącą) sekwencję elementów.
        \item \textit{Postorder (lewe, prawe, korzeń):} Wykorzystywany przy bezpiecznym zwalnianiu pamięci węzłów.
    \end{itemize}
    
    \item \textbf{Usuwanie elementu (Deletion):}
    Jest to najbardziej złożona operacja, wymagająca rozpatrzenia trzech odrębnych przypadków:
    \begin{enumerate}
        \item \textit{Usuwanie liścia (węzeł bez potomków):} Wystarczy usunąć węzeł z pamięci i ustawić wskaźnik rodzica na \texttt{nullptr}.
        \item \textit{Usuwanie węzła z jednym dzieckiem:} Wskaźnik z rodzica usuwanego węzła zostaje przekierowany bezpośrednio na jego jedyne dziecko, po czym usuwany węzeł jest zwalniany.
        \item \textit{Usuwanie węzła z dwojgiem dzieci:} Należy odnaleźć bezpośrednika usuwanego węzła -- zazwyczaj jest to najmniejszy element w prawym poddrzewie (tzw. następnik, ang. \textit{inorder successor}). Wartość następnika jest kopiowana do usuwanego węzła, a następnie rekurencyjnie usuwa się następnik (który ma co najwyżej jedno dziecko).
    \end{enumerate}
\end{itemize}

\subsection{Przykładowy przebieg działania algorytmu krok po kroku}

\hspace{0.60cm}
W celu zaprezentowania poprawnego działania zaimplementowanych algorytmów prześledzono proces tworzenia drzewa BST oraz wykonania na nim serii operacji dla konkretnego zestawu danych wejściowych.

\hspace{0.60cm}
Zestaw danych wejściowych stanowi ciąg 7 liczb całkowitych przekazany do pustego drzewa w następującej kolejności:
\begin{equation}
    S = \{50, 30, 70, 20, 40, 60, 80\}
\end{equation}

\hspace{0.60cm}
Przebieg budowy drzewa BST krok po kroku przedstawia tabela \ref{tab:przebieg_budowy} (s. \pageref{tab:przebieg_budowy}).

\begin{tabela}
	{Analiza budowy drzewa BST krok po kroku}
	{Analiza budowy drzewa BST krok po kroku}
	{
		\begin{tabular}{|c|c|p{7cm}|} \hline
			\textbf{Krok} & \textbf{Wartość} & \textbf{Stan struktury i podjęte decyzje algorytmiczne} \\ \hline
			1 & 50 & Drzewo puste. Utworzenie korzenia z wartością 50. \\ \hline
			2 & 30 & $30 < 50 \rightarrow$ wstawienie jako lewe dziecko węzła 50. \\ \hline
			3 & 70 & $70 > 50 \rightarrow$ wstawienie jako prawe dziecko węzła 50. \\ \hline
			4 & 20 & $20 < 50 \rightarrow$ lewo; $20 < 30 \rightarrow$ lewe dziecko węzła 30. \\ \hline
			5 & 40 & $40 < 50 \rightarrow$ lewo; $40 > 30 \rightarrow$ prawe dziecko węzła 30. \\ \hline
			6 & 60 & $60 > 50 \rightarrow$ prawo; $60 < 70 \rightarrow$ lewe dziecko węzła 70. \\ \hline
			7 & 80 & $80 > 50 \rightarrow$ prawo; $80 > 70 \rightarrow$ prawe dziecko węzła 70. \\ \hline
		\end{tabular}
	}
	\label{tab:przebieg_budowy}
\end{tabela}

\hspace{0.60cm}
Po wstawieniu wszystkich elementów otrzymujemy drzewo o wysokości $h = 2$, którego stan tekstowy w poszczególnych metodach przechodzenia prezentuje się następująco:
\begin{itemize}
    \item \textbf{Preorder:} $50 \rightarrow 30 \rightarrow 20 \rightarrow 40 \rightarrow 70 \rightarrow 60 \rightarrow 80$
    \item \textbf{Inorder:} $20 \rightarrow 30 \rightarrow 40 \rightarrow 50 \rightarrow 60 \rightarrow 70 \rightarrow 80$ (uporządkowane)
    \item \textbf{Postorder:} $20 \rightarrow 40 \rightarrow 30 \rightarrow 60 \rightarrow 80 \rightarrow 70 \rightarrow 50$
\end{itemize}

\hspace{0.60cm}
Kolejnym krokiem w przykładzie jest \textbf{wyszukanie drogi do elementu 40}:
\begin{enumerate}
    \item Start w korzeniu (wartość 50). Dodanie 50 do ścieżki.
    \item Szukana wartość $40 < 50$, przechodzimy w lewo do węzła 30. Dodanie 30 do ścieżki.
    \item Szukana wartość $40 > 30$, przechodzimy w prawo do węzła 40. Dodanie 40 do ścieżki.
    \item Odnaleziono szukaną wartość $40 == 40$.
    \item Zwrócona ścieżka: \texttt{[50, 30, 40]}.
\end{enumerate}

\hspace{0.60cm}
Następnie wykonujemy operację \textbf{usunięcia węzła 30} (węzeł posiadający dwoje dzieci: 20 oraz 40):
\begin{enumerate}
    \item Zlokalizowanie węzła 30 w lewym poddrzewie korzenia.
    \item Odnalezienie następnika w prawym poddrzewie węzła 30. Najmniejsza wartość w prawym poddrzewie węzła 30 to węzeł 40.
    \item Skopiowanie wartości 40 do węzła 30.
    \item Usunięcie oryginalnego węzła 40 (który był liściem).
\end{enumerate}

\hspace{0.60cm}
Stan drzewa po usunięciu elementu 30 w przejściu \textit{Inorder} wynosi:
\begin{equation}
    \text{Inorder}_{\text{po usunięciu}} = \{20, 40, 50, 60, 70, 80\}
\end{equation}

\hspace{0.60cm}
Zimplementowaną logikę usuwania węzła z dwójką dzieci przedstawia poniższy fragment kodu w języku C++ (listing \ref{lst:listing_remove}).

\begin{lstlisting}[caption=Algorytm usuwania węzła z dwójką dzieci w klasie BSTTree, label={lst:listing_remove}, language=C++]
Node* BSTTree::removeNode(Node* node, int key) {
    if (node == nullptr) return nullptr;

    if (key < node->value) {
        node->left = removeNode(node->left, key);
    } else if (key > node->value) {
        node->right = removeNode(node->right, key);
    } else {
        // Wezel znaleziony
        if (node->left == nullptr) {
            Node* temp = node->right;
            delete node;
            return temp;
        } else if (node->right == nullptr) {
            Node* temp = node->left;
            delete node;
            return temp;
        }
        // Case: Dwoje dzieci - szukanie nastepnika (minValueNode)
        Node* temp = findMin(node->right);
        node->value = temp->value;
        node->right = removeNode(node->right, temp->value);
    }
    return node;
}
\end{lstlisting}

\subsection{Opis narzędzi: System kontroli wersji Git oraz platforma GitHub}

\hspace{0.60cm}
Współczesna inżynieria oprogramowania wymaga stosowania zaawansowanych narzędzi wspomagających pracę grupową, automatyzację testów oraz kontrolę jakości kodu. W realizowanym projekcie wykorzystano następujący zestaw narzędzi programistycznych:

\begin{enumerate}
    \item \textbf{Git (Rozproszony System Kontroli Wersji):}
    Git śledzi historię zmian w plikach projektowych. Każda zmiana jest rejestrowana w postaci bezpiecznego commita opatrzonego unikalnym sumarycznym kodem skrótu SHA-1. Kluczowe pojęcia architektoniczne systemu Git obejmują:
    \begin{itemize}
        \item \textit{Working Directory (Katalog roboczy):} Pliki na dysku, na których programista bezpośrednio wykonuje pracę.
        \item \textit{Staging Area / Index (Poczekalnia):} Obszar przygotowawczy, do którego za pomocą polecenia \texttt{git add} trafiają pliki przeznaczone do najnowszego commita.
        \item \textit{Local Repository (Lokalne repozytorium):} Baza danych na lokalnym dysku przechowująca pełną historię wersji po wykonaniu polecenia \texttt{git commit}.
        \item \textit{Remote Repository (Zdalne repozytorium):} Centralna kopia projektu na serwerze (GitHub).
    \end{itemize}

    \item \textbf{Platforma GitHub i praca na gałęziach (Branching):}
    GitHub umożliwia udostępnianie kodu, prowadzenie dyskusji oraz integrację prac programistów. Izolacja wprowadzanych funkcji realizowana jest poprzez gałęzie (\textit{branches}). Praca w projekcie dwuosobowym wymagała opanowania mechanizmu scalania (\textit{git merge}) oraz radzenia sobie z konfliktami.
    
    Konflikt powstaje, gdy dwie osoby zmodyfikują te same linie w tym samym pliku na odrębnych gałęziach. Wtedy system Git zatrzymuje proces automatycznego scalania i nakłada znaczniki konfliktu:
\begin{verbatim}
<<<<<<< HEAD
    // Kod rozwiniety przez Osobe A na galezi glownej
=======
    // Kod zmodyfikowany przez Osobe B na galezi roboczej
>>>>>>> feature-branch
\end{verbatim}
    Rozwiązanie konfliktu polega na ręcznej analizie zmian, wyborze właściwej wersji kodu i ponownym wykonaniu zatwierdzenia commitem.

    \item \textbf{Integracja ze środowiskiem IDE (Visual Studio / VS Code):}
    Do zarzadzania repozytorium wykorzystano wbudowaną wtyczkę Git w środowiskach Microsoft Visual Studio / Visual Studio Code. Umożliwia ona wizualne porównywanie różnic w kodzie (ang. \textit{diff}), szybkie rozwiązywanie konfliktów przyciskami ,,Accept Current Change'' / ,,Accept Incoming Change'' oraz graficzne przeglądanie drzewa historii commitów.

    \item \textbf{Generator dokumentacji Doxygen:}
    Doxygen jest narzędziem służącym do automatycznego tworzenia dokumentacji technicznej w formatach HTML i \LaTeX{} bezpośrednio z komentarzy umieszczonych w kodzie źródłowym C++. Użycie specjalnych tagów (np. \texttt{@brief}, \texttt{@param}, \texttt{@return}) pozwala uzyskać przejrzystą, czytelną dla innych programistów specyfikację klas.
\end{enumerate}

\hspace{0.60cm}
Podsumowując analizę problemu, omówione algorytmy i zasady operowania na strukturze BST oraz poznane narzędzia kontroli wersji stanowią fundament do dalszego opisu implementacji i wyników testów w kolejnych rozdziałach pracy.